% Generated draft: requires mathematical and source review.
\documentclass[11pt]{article}
\usepackage[margin=25mm]{geometry}
\usepackage{fontspec}
\setmainfont{Latin Modern Roman}
\usepackage{amsmath,amssymb,mathtools,unicode-math}
\setmathfont{Latin Modern Math}
\usepackage{enumitem,needspace,xurl,hyperref}
\setlength{\parindent}{0pt}
\setlength{\parskip}{5pt}
\setlength{\emergencystretch}{3em}
\newcommand{\N}{\mathbb{N}}
\newcommand{\Z}{\mathbb{Z}}
\newcommand{\Q}{\mathbb{Q}}
\newcommand{\R}{\mathbb{R}}
\newcommand{\C}{\mathbb{C}}
\newcommand{\F}{\mathbb{F}}
\newcommand{\A}{\mathbb{A}}
\newcommand{\PP}{\mathbb{P}}
\newcommand{\E}{\mathbb{E}}
\newcommand{\eps}{\varepsilon}
\newcommand{\dd}{\,\mathrm{d}}
\begin{document}
\section*{UnsolvedMath ID: 30000279}
\textbf{Source problem code:} catalogue entry
\subsection{Definitions and mathematical statement}
Let $G$ be an arbitrary Carnot group: a connected simply connected nilpotent Lie group whose Lie algebra has a stratification $\mathfrak g=V_1\oplus\cdots\oplus V_s$ with $[V_1,V_j]=V_{j+1}$ and $V_{s+1}=0$. Let $X_1,\ldots,X_m$ be left-invariant vector fields from a basis of the horizontal layer $V_1$, and let $\dd g$ denote Haar measure. In the usual finite-valued convention, an $H$-convex function is convex on horizontal line segments $t\mapsto g\exp(tv)$, $v\in V_1$. Its derivatives in this question are distributional, using the locally integrable representative, rather than assuming $C^2$ regularity.

The conjecture asks whether, for every such group and every $H$-convex $u$, each distribution
\[
 [X_i,X_j]u=(X_iX_j-X_jX_i)u
\]
is represented by a signed Radon measure, for all $1\le i,j\le m$. Since these left-invariant fields are divergence-free for Haar measure, this means that there exists $\mu_{ij}$ with locally finite total variation such that
\[
 -\int_Gu\,[X_i,X_j]\varphi\,\dd g
       =\int_G\varphi\,\dd\mu_{ij}
 \qquad(\varphi\in C_c^\infty(G)).
\]
Equivalently, for each compact $K\subset G$ there should be a constant $C_K$ such that
\[
 \left|\int_Gu\,[X_i,X_j]\varphi\,\dd g\right|
 \le C_K\|\varphi\|_\infty
 \quad\text{whenever }\operatorname{supp}\varphi\subset K.
\]
The measure need not be nonnegative: the commutator is antisymmetric in $i,j$. Radon means locally finite variation, not finite mass on the whole noncompact group. No restriction to step two or to Heisenberg groups is imposed, and no classical second differentiability is added.

The literature assessment checked on 21 August 2026 reports positive results in Heisenberg and several step-two settings but does not cover arbitrary Carnot groups. Its open-status claim is supplied dated triage requiring independent verification.
\subsection{Short English statement}
Do all horizontal commutator derivatives of an $H$-convex function on any Carnot group define locally finite signed measures? The question is distributional and includes groups of arbitrary step, beyond the special cases in the supplied assessment.
\subsection{Sources}
\begin{itemize}
\item \emph{Partielle Differentialgleichungen} (2006), Oberwolfach Reports, DOI 10.4171/owr/2005/33, pp. 1867--1928, \url{https://doi.org/10.4171/owr/2005/33}; proposed year supplied as 2005.
\item Literature assessment checked 21 August 2026: original report, \url{https://doi.org/10.4171/owr/2005/33}; ``Carnot-group regularity survey,'' \url{https://people.dm.unipi.it/magnani/works/Ober05.pdf}.
\item Local UnsolvedMath registry, permanent ID $30000279$, exact source code \texttt{catalogue entry}, created 25 August 2026 and updated 16 September 2026. Supplied verification restores notation only and preserves all quantifiers.
\end{itemize}
\end{document}
